\chapter{Validación de estructura}
\label{ap:estructura}

\section*{Dimensiones del tanque}

En base a la relación estandarizada de los biorreactores:

\begin{align}
A_r &= \pi(107.5mm)^2 = 36\,305.115\ \text{mm}^2 \\
A_d &= \pi(48.26mm)^2 = 7\,316.87\ \text{mm}^2 \\
\frac{A_d}{A_r} &= \frac{7\,316.87}{36\,305.115} = 0.201
\end{align}


La Tabla~\ref{fig:radio_curvatura} muestra las relaciones del tamaño del radio de curvatura del tubo exterior
de retorno.
\begin{figure}[H]
	\centering
	\includegraphics[width=\textwidth]{relaciones_radio_curvatura}
	\caption{Tabla de relaciones del tamaño del radio para el tubo exterior de retorno.}
	\label{fig:radio_curvatura}
\end{figure}



\section*{Validación de pernos}

Para la validación de selección de los tornillos se analizó la conexión de un cilindro
de acero con un recipiente a presión de acero inoxidable AISI~304 mediante 8~pernos. Un sello de junta
confinado tiene un diámetro de sellado efectivo de 170~mm. El cilindro almacena
gas a una presión máxima de 5~bar.

\begin{figure}[H]
  \centering
  \includegraphics[width=\textwidth]{tornillos_rep_visual}
  \caption{Representación visual del problema de los tornillos.}
  \label{fig:tornillos}
\end{figure}

Datos conocidos:
\begin{itemize}
  \item Di\'ametro sellado efectivo: $d_{sel} = 170$~mm, espesor tapa: $t = 5$~mm
  \item $N = 8$~pernos, $P_g = 5$~bar ($0.50$~MPa), especificaci\'on: 5/16-18~UNC~$\times$~1\textonequarter~in, SAE Grado~5
\end{itemize}

Los c\'alculos completos del modelo de rigidez de Shigley se presentan a
continuaci\'on, con las propiedades del perno 5/16-18~UNC en pulgadas:

\begin{align}
d &= 0.3125\ \text{in}, \quad A_t = 0.05219\ \text{in}^2, \quad A_d = 0.07670\ \text{in}^2 \\
l_d &= 0.375\ \text{in}, \quad l_t = 0.875\ \text{in}, \quad E = 30{\times}10^6\ \text{psi}
\end{align}

\begin{equation}
k_b = \frac{A_d A_t E}{A_d l_t + A_t l_d} = 3.106 \times 10^6\ \frac{\text{lbf}}{\text{in}}
\end{equation}

\begin{equation}
k_m = 11.085 \times 10^6\ \frac{\text{lbf}}{\text{in}}, \quad
C = \frac{k_b}{k_b + k_m} = \frac{3.106}{3.106 + 11.085} = 0.2188
\end{equation}

\begin{equation}
S_p = 85\,000\ \text{psi} \quad \text{(SAE Grado~5)}, \quad
F_i = 0.75\, A_t\, S_p = 3{,}340.5\ \text{lbf}
\end{equation}

\begin{equation}
P_{total} = 0.50\ \text{MPa} \times \pi\left(\frac{170}{2}\right)^2\ \text{mm}^2
= 5{,}004.8\ \text{lbf}, \quad P = \frac{P_{total}}{8} = 625.6\ \text{lbf/perno}
\end{equation}

\begin{align}
n_p &= \frac{S_p A_t}{C P + F_i} = \frac{85{,}000 \times 0.05219}{0.2188 \times 625.6 + 3{,}340.5} = 1.281 \\
n_L &= \frac{S_p A_t - F_i}{C P} = \frac{85{,}000 \times 0.05219 - 3{,}340.5}{0.2188 \times 625.6} = 8.133 \\
n_0 &= \frac{F_i}{P(1-C)} = \frac{3{,}340.5}{625.6 \times (1 - 0.2188)} = 6.836
\end{align}

Los tres factores de seguridad ($n_p = 1.281$, $n_L = 8.133$, $n_0 = 6.836$)
superan la unidad, confirmando que 8~tornillos 5/16-18~UNC~$\times$~1\textonequarter~in
SAE Grado~5 resisten adecuadamente la aplicaci\'on.

\section*{Presión hidrostática}

Para corroborar que el biorreactor soporta el volumen de agua ingresado se
sometió a un estudio de tensión estática. Se aplicaron condiciones de presión
hidrostática (densidad del agua: $9.170 \times 10^{-7}$~kg/mm$^3$) y una carga de gravedad
estándar de $9.807$~m/s$^2$.

\begin{figure}[H]
  \centering
  \includegraphics[width=\textwidth]{fea_hidro_cargas}
  \caption{Aplicación de cargas a las caras internas del biorreactor.}
  \label{fig:cargas_hidrostaticas}
\end{figure}

Con los resultados a través de la simulación, se comprobó que la estructura del
biorreactor soporta con un gran margen la cantidad de líquido ingresado, La tensi\'on de von Mises m\'axima fue de $2.36$~MPa,
resultando en un factor de seguridad m\'inimo de $\mathbf{91.1}$, amplio margen
frente a cualquier condici\'on de operaci\'on normal.

\begin{figure}[H]
  \centering
  \includegraphics[width=\textwidth]{fea_hidro_fs}
  \caption{Resultados del factor de seguridad al aplicar una presión hidrostática.}
  \label{fig:fs_hidrostatica}
\end{figure}

\begin{figure}[H]
  \centering
  \includegraphics[width=\textwidth]{fea_hidro_vonmises}
  \caption{Mapa de esfuerzo de von Mises --- presión hidrostática (máx = 2.36~MPa).}
  \label{fig:hidro_vonmises}
\end{figure}

\section*{Presión manométrica}

El biorreactor opera con una presi\'on manom\'etrica nominal de 2~bar (0.20~MPa).
Para la simulaci\'on se aplic\'o una carga conservadora de $\mathbf{0.50~MPa}$ (5~bar)
a todas las caras internas, representando el peor caso esperado.

\begin{figure}[H]
  \centering
  \includegraphics[width=\textwidth]{fea_mano_cargas}
  \caption{Aplicación de cargas de 0.50~MPa (5~bar) a las caras internas del biorreactor.}
  \label{fig:cargas_manometricas}
\end{figure}

Los resultados obtenidos mostraron que la estructura soporta esta presi\'on extrema
con una tensi\'on de von Mises m\'axima de $67.7$~MPa y un factor de seguridad
m\'inimo de $\mathbf{3.18}$, superando el criterio de dise\'no FS~$\geq 2$. El modelo fue analizado
utilizando una malla de elementos sólidos parabólicos con un total de 121\,510
nodos y 61\,028 elementos.

\begin{figure}[H]
  \centering
  \includegraphics[width=\textwidth]{fea_mano_fs}
  \caption{Resultados del factor de seguridad al aplicar una presión de 0.50~MPa o 5~bar (FS~mín~=~3.18).}
  \label{fig:fs_manometrica}
\end{figure}

\begin{figure}[H]
  \centering
  \includegraphics[width=\textwidth]{fea_mano_vonmises}
  \caption{Mapa de esfuerzo de von Mises --- presión manométrica 0.50~MPa (máx = 67.7~MPa).}
  \label{fig:mano_vonmises}
\end{figure}

\begin{figure}[H]
  \centering
  \includegraphics[width=\textwidth]{fea_mano_desp}
  \caption{Mapa de desplazamiento total --- presión manométrica 0.50~MPa.}
  \label{fig:mano_desp}
\end{figure}

\section*{Diseño para tolerancias de agujeros pasados y roscados}

Para garantizar las dimensiones de los agujeros pasados para la tornillería de 5/16, 1/2, 5/8 y 3/4~de pulgada se utilizaron las especificaciones según la norma ASME~B18.2.8. De acuerdo con la Tabla~2 del estándar para un ajuste normal (\textit{Normal Fit}), se estableció el uso de brocas nominales de 11/32, 9/16, 11/16 y 13/16~de pulgada. En los planos de fabricación, los diámetros se acotaron utilizando los tamaños mínimos de la herramienta correspondientes a 0.344, 0.562, 0.688 y 0.812~pulgadas. Para asegurar que el mecanizado se mantenga dentro de la norma, se aplicó una tolerancia de $+0.010$~pulgadas para los barrenos de 5/16'', 1/2'' y 5/8'', mientras que para el barreno de 3/4'' se utilizó una tolerancia de $+0.012$~pulgadas; esto garantiza que en ningún caso se excedan los límites máximos permisibles de 0.354, 0.572, 0.698 y 0.824~pulgadas dictados por el estándar~\cite{ASME_B18_2_8}. Mientras que para las dimensiones de los agujeros pasados para los elementos de fijación métricos M5, se realizó con la Tabla~3 del estándar para un ajuste normal (\textit{Normal Fit}), donde se estableció el uso de una broca nominal de 5.5~mm~\cite{ASME_B18_2_8}.

\begin{figure}[H]
  \centering
  \includegraphics[width=0.9\textwidth]{tabla_2}
  \caption{Tabla~2 de la norma ASME~B18.2.8: agujeros de paso para tornillería en pulgadas~\cite{ASME_B18_2_8}.}
  \label{fig:tabla_2}
\end{figure}

\begin{figure}[H]
  \centering
  \includegraphics[width=0.9\textwidth]{tabla_3}
  \caption{Tabla~3 de la norma ASME~B18.2.8: agujeros de paso para elementos de fijación métricos~\cite{ASME_B18_2_8}.}
  \label{fig:tabla_3}
\end{figure}

Para garantizar el sellado hermético en las conexiones, las tolerancias de las roscas cónicas para tubería de 1/4 y 3/4~de pulgada se definieron utilizando las especificaciones de la norma ASME~B1.20.1. De acuerdo con el estándar para roscas NPT (\textit{National Pipe Taper}), se estableció el uso de 18~hilos por pulgada (TPI) para la conexión de 1/4'' y de 14~hilos por pulgada para la de 3/4''. Esto garantiza que la variación máxima en el acoplamiento no exceda el límite permisible de $\pm 1$~vuelta de rosca respecto al plano de apriete manual (\textit{hand-tight plane}), asegurando la interferencia necesaria en los flancos para lograr el sellado dictado por la norma~\cite{ASME_B1_20_1}.

\section*{Diseño de ranura para O-ring}

Para garantizar el sellado hermético mediante juntas tóricas (O-rings) con secciones transversales nominales de 3/16 y 1/4~de pulgada, se utilizaron las especificaciones de diseño según el Parker O-Ring Handbook. De acuerdo con las tablas del estándar para aplicaciones de sellado estático, se estableció la geometría de las ranuras (\textit{glands}) (G) considerando las secciones transversales reales de los anillos de $0.280 \pm 0.010$ y $0.352 \pm 0.010$~pulgadas. En los planos de fabricación, la profundidad (L) de la ranura se acotó utilizando la dimensión recomendada de $0.157 \pm 0.05$~pulgadas para el O-ring de 3/16'' y de $0.206 \pm 0.005$~pulgadas para el de 1/4''~\cite{parker_oring}.

\begin{figure}[H]
  \centering
  \includegraphics[width=0.9\textwidth]{tabla_orings}
  \caption{Carta de diseño de ranura para O-ring tipo sello de cara (Parker O-Ring Handbook)~\cite{parker_oring}.}
  \label{fig:tabla_orings}
\end{figure}

\section*{Vistas explosionadas}

% ── Ensamble biorreactor (7 páginas) ───────────────────────────────────────
\begin{figure}[H]
  \centering
  \includegraphics[width=\textwidth,page=1]{ensamble_biorreactor}
  \caption{Vista explosionada del biorreactor completo.}
  \label{fig:ensamble_biorreactor_completo}
\end{figure}

\begin{figure}[H]
  \centering
  \includegraphics[width=\textwidth,page=2]{ensamble_biorreactor}
  \caption{Vista explosionada del biorreactor, parte 1.}
  \label{fig:ensamble_biorreactor_p1}
\end{figure}

\begin{figure}[H]
  \centering
  \includegraphics[width=\textwidth,page=3]{ensamble_biorreactor}
  \caption{Vista explosionada del biorreactor, parte 2.}
  \label{fig:ensamble_biorreactor_p2}
\end{figure}

\begin{figure}[H]
  \centering
  \includegraphics[width=\textwidth,page=4]{ensamble_biorreactor}
  \caption{Vista explosionada del biorreactor, parte 3.}
  \label{fig:ensamble_biorreactor_p3}
\end{figure}

\begin{figure}[H]
  \centering
  \includegraphics[width=\textwidth,page=5]{ensamble_biorreactor}
  \caption{Vista explosionada del biorreactor, parte 4.}
  \label{fig:ensamble_biorreactor_p4}
\end{figure}

\begin{figure}[H]
  \centering
  \includegraphics[width=\textwidth,page=6]{ensamble_biorreactor}
  \caption{Vista explosionada del biorreactor, parte 5.}
  \label{fig:ensamble_biorreactor_p5}
\end{figure}

\begin{figure}[H]
  \centering
  \includegraphics[width=\textwidth,page=7]{ensamble_biorreactor}
  \caption{Vista explosionada del biorreactor, parte 6.}
  \label{fig:ensamble_biorreactor_p6}
\end{figure}

% ── Ensamble HMI (6 páginas) ───────────────────────────────────────────────
\begin{figure}[H]
  \centering
  \includegraphics[width=\textwidth,page=1]{ensamble_hmi}
  \caption{Vista explosionada de la HMI completa.}
  \label{fig:ensamble_hmi_completa}
\end{figure}

\begin{figure}[H]
  \centering
  \includegraphics[width=\textwidth,page=2]{ensamble_hmi}
  \caption{Vista explosionada de la HMI, parte 1.}
  \label{fig:ensamble_hmi_p1}
\end{figure}

\begin{figure}[H]
  \centering
  \includegraphics[width=\textwidth,page=3]{ensamble_hmi}
  \caption{Vista explosionada de la HMI, parte 2.}
  \label{fig:ensamble_hmi_p2}
\end{figure}

\begin{figure}[H]
  \centering
  \includegraphics[width=\textwidth,page=4]{ensamble_hmi}
  \caption{Vista explosionada de la HMI, parte 3.}
  \label{fig:ensamble_hmi_p3}
\end{figure}

\begin{figure}[H]
  \centering
  \includegraphics[width=\textwidth,page=5]{ensamble_hmi}
  \caption{Vista explosionada de la HMI, parte 4.}
  \label{fig:ensamble_hmi_p4}
\end{figure}

\begin{figure}[H]
  \centering
  \includegraphics[width=\textwidth,page=6]{ensamble_hmi}
  \caption{Vista explosionada de la HMI, parte 5.}
  \label{fig:ensamble_hmi_p5}
\end{figure}

% ── Ensamble bombas (1 página) ─────────────────────────────────────────────
\begin{figure}[H]
  \centering
  \includegraphics[width=\textwidth,page=1]{ensamble_bombas}
  \caption{Vista explosionada de las bombas.}
  \label{fig:ensamble_bombas}
\end{figure}
